% Ekdant v0.1.0 -- progress report after the supervisor meeting of 2026-11-18
% Build: tectonic Ekdant_v0.1.0_Report.tex   (XeLaTeX engine)
% Sources of every number: 04-Model-DeepDive/30_L2_v6_Roles_Spec.md (design,
% measured size), run.md A3t (protocol), 07-Management/10_Known_Issues.md T64/T65,
% 09-Code (code, configs, tests: 404 passed), the 2026-11-13 progress report
% (v0.0.x results).  v0.1.0 has not been trained: no performance number of it
% appears here.
\documentclass[a4paper,12pt,parskip=half,headsepline]{scrartcl}

\usepackage{amsmath}
\usepackage{fontspec}
% TeX Gyre Pagella (text and maths) and Heros (headings), loaded by file name
\setmainfont{texgyrepagella}[Extension=.otf,UprightFont=*-regular,
  BoldFont=*-bold,ItalicFont=*-italic,BoldItalicFont=*-bolditalic]
\setsansfont{texgyreheros}[Extension=.otf,UprightFont=*-regular,
  BoldFont=*-bold,ItalicFont=*-italic,BoldItalicFont=*-bolditalic]
\usepackage{unicode-math}
\setmathfont{texgyrepagella-math.otf}
\linespread{1.08}
\usepackage[a4paper,margin=2.5cm,top=2.8cm,bottom=2.8cm,headsep=0.6cm]{geometry}
\usepackage{graphicx}
\usepackage{float}
\usepackage{booktabs,tabularx,array,multirow}
\usepackage{xcolor}
\usepackage{tikz}
\usetikzlibrary{arrows.meta,positioning,calc,fit,backgrounds,shapes.geometric}
\usepackage[font=small,labelfont=bf,margin=0.5cm]{caption}
\usepackage{enumitem}
\usepackage[most]{tcolorbox}
\usepackage{microtype}
\usepackage[hidelinks]{hyperref}
\usepackage{scrlayer-scrpage}

\definecolor{lblue}{HTML}{2A78D6}
\definecolor{lsoft}{HTML}{E6F0FB}
\definecolor{lorange}{HTML}{EB6834}
\definecolor{lgreen}{HTML}{1F7A46}
\definecolor{ink}{HTML}{17191C}
\definecolor{muted}{HTML}{575B62}
\definecolor{rule}{HTML}{DEDFDC}
\definecolor{rA}{HTML}{2A78D6}
\definecolor{rB}{HTML}{7B4FB8}
\definecolor{rC}{HTML}{C2185B}
\definecolor{rD}{HTML}{1F8A80}
\definecolor{plainbg}{HTML}{F3F7EE}
\definecolor{plainfr}{HTML}{1F7A46}
\definecolor{techbg}{HTML}{F4F4F6}

\addtokomafont{disposition}{\color{ink}}
\addtokomafont{section}{\Large}
\addtokomafont{subsection}{\large}
\addtokomafont{caption}{\small}
\addtokomafont{captionlabel}{\sffamily\bfseries}
\setkomafont{pageheadfoot}{\normalfont\sffamily\small\color{muted}}
\setkomafont{pagenumber}{\normalfont\sffamily\small}
\setlist{itemsep=2pt,topsep=3pt}
\renewcommand{\arraystretch}{1.18}
\newcolumntype{Y}{>{\raggedright\arraybackslash}X}

\newcommand{\sig}{\ensuremath{^{\ast}}}
\newcommand{\code}[1]{\texttt{\small #1}}

\newtcolorbox{keybox}[1][]{enhanced,colback=lsoft,colframe=lblue,boxrule=0.6pt,
  arc=2pt,left=8pt,right=8pt,top=6pt,bottom=6pt,fonttitle=\bfseries,#1}
\newtcolorbox{rulebox}[1][]{enhanced,colback=white,colframe=ink,boxrule=0.5pt,
  arc=2pt,left=8pt,right=8pt,top=6pt,bottom=6pt,fonttitle=\bfseries,
  colbacktitle=ink!8,coltitle=ink,#1}
% the two ways of explaining every part (author's request)
\newtcolorbox{plain}[1][]{enhanced,breakable,colback=plainbg,colframe=plainfr,boxrule=0.5pt,
  arc=2pt,left=8pt,right=8pt,top=5pt,bottom=5pt,fonttitle=\bfseries\small,
  title={In plain words},fonttitle=\sffamily\bfseries\small,
  before upper={\setlength{\parskip}{5pt}},#1}
\newtcolorbox{tech}[1][]{enhanced,breakable,colback=techbg,colframe=muted,boxrule=0.5pt,
  arc=2pt,left=8pt,right=8pt,top=5pt,bottom=5pt,fonttitle=\bfseries\small,
  title={Technical explanation},fonttitle=\sffamily\bfseries\small,
  before upper={\setlength{\parskip}{5pt}},#1}

\pagestyle{scrheadings}
\clearpairofpagestyles
\ihead{Ekdant v0.1.0 \textbullet{} the sandbox with four role agents}
\ohead{H.~Gorasiya \textbullet{} 18 November 2026}
\cfoot{\pagemark}

\tikzset{
  >={Stealth[length=2.4mm]},
  blk/.style={draw=ink,thick,rounded corners=3pt,fill=white,align=center,
              minimum height=1.0cm,inner sep=4pt,font=\small},
  sbx/.style={blk,draw=lblue,fill=lsoft,very thick},
  role/.style={draw=#1,fill=#1!10,thick,rounded corners=3pt,align=center,
               minimum width=3.4cm,minimum height=0.95cm,font=\small,inner sep=3pt},
  cand/.style={draw=#1,fill=#1!15,thick,circle,minimum size=0.62cm,inner sep=0pt,font=\footnotesize},
  lab/.style={font=\footnotesize,color=muted,align=center},
  flow/.style={->,thick,ink},
  bflow/.style={->,very thick,lblue},
  thin/.style={->,semithick,muted},
}

\hypersetup{pdftitle={Ekdant v0.1.0: the sandbox with four role agents},
  pdfauthor={Harshil Gorasiya}}

\begin{document}

% =====================================================================
\begin{titlepage}
  \thispagestyle{empty}
  \vspace*{0.6cm}
  {\color{muted}\large\sffamily MSc thesis \textbullet{} Hochschule Heilbronn}\par
  \vspace{0.6cm}
  {\color{ink}\fontsize{25}{30}\selectfont\sffamily\bfseries Ekdant v0.1.0:\\[3pt] a sandbox of four role agents\\[3pt] inside a small language model\par}
  \vspace{0.6cm}
  {\Large\sffamily\color{lblue} Architecture, new evaluation and status after the meeting of 18 November 2026\par}
  \vspace{0.7cm}
  \begin{tikzpicture}\draw[lblue,line width=2pt] (0,0) -- (14,0);\end{tikzpicture}\par
  \vspace{0.5cm}
  {\large\sffamily Harshil Gorasiya\par}
  \vspace{0.2cm}
  {\sffamily\color{muted} Status: 18 November 2026\par}
  \vfill
  \begin{keybox}[title={In one paragraph}]
  \small
  The first sandbox (v0.0.1--v0.0.5) placed four copies of each word's state at
  \emph{random} positions, let them debate and averaged them; it never improved
  reasoning, and its best gain was a preference for answering ``yes''.
  Following the meeting, \textbf{Ekdant v0.1.0} rebuilds the sandbox as a small team
  of agents \emph{inside} the model: the model's own state is kept, four reviewers
  with \textbf{different jobs and rules} (explore, check consistency, express
  uncertainty, build the best-supported alternative) each propose a corrected
  state, they read each other once, and a \textbf{learned judge} decides word by
  word whose version to use --- including ``keep the original''. Each reviewer can
  also get its own training signal. The evaluation moves from multiple choice to
  \textbf{text generation}: the model trains on written medical reasoning and is
  judged on whether its answers contain the key medical facts (MedLFQA).
  \textbf{Status:} built, tested (404 automated tests) and checked end-to-end on
  CPU; \textbf{not trained yet} (no GPU access at the moment), so this report
  contains no performance result of v0.1.0.
  \end{keybox}
  \vspace{0.3cm}
  {\footnotesize\color{muted} Every section explains each part twice: a green box \emph{in plain words}, and a grey box with the \emph{technical explanation}. v0.0.x numbers are means over seeds 17, 23, 31 with a paired bootstrap 95\,\% CI; \sig{} = interval excludes 0.\par}
\end{titlepage}

\tableofcontents

\begin{rulebox}[title={What is new in v0.1.0, at a glance}]
\begin{tabularx}{\linewidth}{@{}p{3.6cm}Y@{}}
\textbf{Roles, not copies} & Four reviewers with different inputs and rules (R1 explore, R2 consistency, R3 uncertainty, R4 best-supported) instead of four random projections.\\[2pt]
\textbf{The original is kept} & The model's own state is candidate 0; the sandbox can decide to change nothing.\\[2pt]
\textbf{A judge, not an average} & A learned gate weighs the five candidates per word; its weights say \emph{which} reviewer shaped the word.\\[2pt]
\textbf{Training per role} & Optional objectives push each reviewer towards its job and stop ignored reviewers from fading out.\\[2pt]
\textbf{Placement is measured} & One insertion point per run, chosen by a flag; a study over depths 5, 10, 15, 20, 25 and the end.\\[2pt]
\textbf{New benchmark} & Text generation: reasoning-text training; MedLFQA scored for comprehensiveness and hallucination.
\end{tabularx}
\end{rulebox}

% =====================================================================
\section{Where we stand and what the meeting changed}

\subsection{The versions}

\begin{center}
\begin{tabularx}{\linewidth}{@{}lY@{}}
\toprule
\textbf{Name} & \textbf{What}\\
\midrule
v0.0.1 & the from-scratch pilot (50\,M parameters, multiple-choice protocol)\\
v0.0.2 -- v0.0.5 & SmolLM2-135M + the first sandbox (``L2''); v0.0.5 = the final model of the 13 November report\\
\textbf{v0.1.0} & \textbf{this report: the redesigned sandbox and the text-generation protocol}\\
v1.0.0 & reserved for the version that is shown to work\\
\bottomrule
\end{tabularx}
\end{center}

\subsection{What v0.0.x showed (13 November report)}

\begin{plain}
The first sandbox did not make the model reason better. Its accuracy gain on
PubMedQA came from answering ``yes'' more often, not from reading the evidence.
When the training data was fixed so that the yes/no answer had to come from the
evidence, the model without any sandbox improved, and the sandbox added nothing on top.
\end{plain}

\begin{center}
\begin{tabularx}{\linewidth}{@{}Ylll@{}}
\toprule
\textbf{Setting (PubMedQA, \% right)} & \textbf{Baseline} & \textbf{v0.0.5 sandbox} & \textbf{Difference}\\
\midrule
Original data & 58.0 & 60.0 & +2.07\sig{} (a ``yes'' preference)\\
Evidence data (T61) & 61.3 & 61.2 & $-0.10$\\
\bottomrule
\end{tabularx}
\end{center}

\subsection{Why the first sandbox could not help --- and the answer of v0.1.0}

\begin{plain}
In the meeting we agreed that the old design had no direction: the four copies
were thrown to random places, nobody had a job, and the result was a plain average.
The model's own understanding was barely used, and the layer could not say why it
chose anything. v0.1.0 keeps the idea of a team working on each word's state, but
gives every member a job and adds a judge.
\end{plain}

\begin{center}
\begin{tabularx}{\linewidth}{@{}YY@{}}
\toprule
\textbf{Problem in v0.0.x} & \textbf{Answer in v0.1.0}\\
\midrule
Four random projections $T_k = W_k h$, all built the same way; no role & Four generators with different inputs and rules (R1--R4, Section~\ref{sec:roles})\\
Output = mean of the four; the model's own state is not a choice & Five candidates including $h$ itself; a learned judge weighs them per word\\
The gate settled at $g \approx 0.027$: the model used about 3\,\% of L2 & The change is sized to the model's stream; the judge, not a fixed gate, decides how much each candidate counts\\
No reason for a choice & The judge's weights per word show which reviewer won (logged, and testable by forcing one candidate)\\
Placed after the last layer, where it can only tilt the final word scores & One configurable insertion point; a placement study over six depths\\
Multiple-choice benchmark rewards a yes/no tilt & Text-generation benchmark that checks the facts in written answers\\
\bottomrule
\end{tabularx}
\end{center}

% =====================================================================
\section{The architecture of the sandbox v0.1.0}

\subsection{The whole model}

\begin{plain}
SmolLM2 reads the text word by word through 30 layers. At \emph{one} chosen point
between these layers, every word has a ``state'' --- 576 numbers holding what the
model currently thinks about it. The sandbox sits exactly there: it takes each
word's state, lets a small team of agents work on it, and hands an improved state
to the remaining layers. Because the rest of the model reads the improved state,
a correction can influence everything that comes after it, including later words.
Nothing else in SmolLM2 changes, and with the sandbox switched off the model is
exactly SmolLM2.
\end{plain}

\begin{figure}[H]
\centering
\resizebox{\linewidth}{!}{%
\begin{tikzpicture}[node distance=0.55cm]
  \node[blk] (txt) {text\\(question)};
  \node[blk,right=of txt] (emb) {embedding};
  \node[blk,right=of emb,minimum width=2.4cm] (low) {decoder layers\\$1 \dots k$};
  \node[sbx,right=1.0cm of low,minimum width=3.0cm,minimum height=1.5cm] (sb) {\textbf{sandbox v0.1.0}\\[-1pt]{\footnotesize h + four role agents}\\[-1pt]{\footnotesize + judge}};
  \node[blk,right=1.0cm of sb,minimum width=2.4cm] (high) {decoder layers\\$k{+}1 \dots 30$};
  \node[blk,right=of high] (norm) {final norm\\+ output head};
  \node[blk,right=of norm] (out) {next\\word};
  \draw[flow] (txt) -- (emb); \draw[flow] (emb) -- (low);
  \draw[bflow] (low) -- node[above,lab]{state $h$} (sb);
  \draw[bflow] (sb) -- node[above,lab]{state $h'$} (high);
  \draw[flow] (high) -- (norm); \draw[flow] (norm) -- (out);
  \node[lab,below=0.35cm of sb] {$k$ is set per run (\code{--addon-layer}):\\after layer 5, 10, 15, 20, 25 or at the end};
  \draw[muted,dashed,rounded corners] ($(low.north west)+(-0.25,0.45)$) rectangle ($(norm.south east)+(0.25,-1.75)$);
  \node[lab,anchor=south west] at ($(low.north west)+(-0.2,0.47)$) {SmolLM2-135M (unchanged)};
\end{tikzpicture}}
\caption{Where the sandbox sits. With the sandbox off (or at step 0 of training) $h' = h$ and the model is exactly SmolLM2.}
\label{fig:whole}
\end{figure}

\begin{tech}
\textbf{Where it is attached.} SmolLM2-135M is a decoder-only transformer with 30
pre-norm decoder layers and hidden size $d = 576$. Each layer adds its result to the
\emph{residual stream} $h \in \mathbb{R}^{B \times T \times d}$ ($B$ sequences of $T$
tokens). The sandbox is inserted with a PyTorch forward hook on decoder layer $k$: the
hook receives the layer's output $h$ and returns $h' = h + \text{Sandbox}(h)$, which is
what layer $k+1$ reads. $k$ is a configuration value (\code{addon\_layer}; $1 \le k \le 30$,
or $k = -1$ for ``after the final norm'', directly before the output head) and can be
overridden per run with \code{--addon-layer}; the transformer's own code is not changed.

\textbf{Why the input is normalised.} The size of the stream differs strongly between
depths: its root-mean-square (rms) is about 2 after the final norm and about 100 inside
the network, with single channels near $2\cdot10^4$. A module that read $h$ directly would
see completely different inputs at different depths. The sandbox therefore reads
$x = \text{RMSNorm}(h)$ (with a learned per-channel weight), computes its change in these
normalised units, and scales the change back by $\text{rms}(h)$ before adding it. The same
weights thus work at every depth, and a change of a given size always means the same
\emph{relative} change of the model's state.

\textbf{Causality.} Every operation is either per token or attends only to earlier tokens
of the same sequence (causal mask, padding excluded). Token $t$ therefore never sees
tokens $> t$, batch rows never mix, and generation with the key--value cache gives
exactly the same output as recomputing the whole sequence (verified by tests at layer 2
and at the end).
\end{tech}

\subsection{Inside the sandbox}

\begin{figure}[H]
\centering
\resizebox{\linewidth}{!}{%
\begin{tikzpicture}
  % input
  \node[blk,minimum width=1.6cm] (h) at (-2.6,-1.9) {word\\state $h$};
  \node[blk,minimum width=2.0cm] (x) at (0,-1.9) {$x = $\\RMSNorm$(h)$};
  \draw[flow] (h) -- (x);
  \node[blk,minimum width=2.4cm,fill=ink!4] (mem) at (0,-4.3) {earlier words\\(causal memory)};
  \draw[thin] (x) -- (mem);
  % roles
  \node[role=rA] (r1) at (5.0,1.3) {\textbf{R1 explore}\\{\footnotesize new direction, $\perp x$}};
  \node[role=rB] (r2) at (5.0,-0.2) {\textbf{R2 consistency}\\{\footnotesize reads earlier words}};
  \node[role=rD] (r4) at (5.0,-1.7) {\textbf{R4 best-supported}\\{\footnotesize words + R1, R2}};
  \node[role=rC] (r3) at (5.0,-3.2) {\textbf{R3 uncertainty}\\{\footnotesize their disagreement}};
  \foreach \r in {r1,r2,r4,r3} \draw[thin] (x.east) -- (\r.west);
  \draw[thin,dashed] (mem.east) to[out=0,in=200] (r2.west);
  \draw[thin,dashed] (mem.east) to[out=0,in=200] (r4.west);
  \draw[thin,rA] (r1.east) to[out=-20,in=20] (r4.east);
  \draw[thin,rB] (r2.east) to[out=-20,in=20] (r4.east);
  \draw[thin,rC] (r4.south) -- node[right,lab]{spread} (r3.north);
  % candidates
  \node[cand=ink] (c0) at (8.6,2.0) {$C_0$};
  \node[cand=rA] (c1) at (8.6,1.0) {$C_1$};
  \node[cand=rB] (c2) at (8.6,0.0) {$C_2$};
  \node[cand=rC] (c3) at (8.6,-1.0) {$C_3$};
  \node[cand=rD] (c4) at (8.6,-2.0) {$C_4$};
  \draw[thin] (x.north) |- (c0.west);
  \draw[thin,rA] (r1.east) -- (c1.west);
  \draw[thin,rB] (r2.east) -- (c2.west);
  \draw[thin,rC] (r3.east) to[out=0,in=180] (c3.west);
  \draw[thin,rD] (r4.east) to[out=0,in=180] (c4.west);
  \node[lab] at (8.6,-2.85) {$C_0 = x$,\\$C_i = x + \Delta_i$};
  % round
  \node[blk,minimum width=2.5cm,minimum height=4.6cm,fill=ink!3] (rnd) at (11.2,0.0) {\textbf{communication}\\\textbf{round}\\[4pt]{\footnotesize each reads}\\{\footnotesize the other four,}\\{\footnotesize adjusts its own}\\{\footnotesize $\Delta_i$; $C_0$ fixed}};
  \foreach \c in {c0,c1,c2,c3,c4} \draw[thin] (\c.east) -- (\c.east -| rnd.west);
  % judge
  \node[sbx,minimum width=2.6cm,minimum height=2.2cm] (gate) at (14.4,0.0) {\textbf{judge}\\{\footnotesize scores $s_0..s_4$}\\{\footnotesize $p = \text{softmax}(s)$}};
  \draw[bflow] (rnd) -- (gate);
  \node[blk,minimum width=3.4cm,minimum height=1.3cm] (out) at (14.4,-3.2) {$h' = h + \text{rms}(h)\sum_{i} p_i \Delta_i$};
  \draw[bflow] (gate) -- (out);
  \node[lab] at (14.75,1.75) {reads each candidate,\\$x$ and the text so far};
\end{tikzpicture}}
\caption{Inside the sandbox v0.1.0. Coloured boxes are the four role generators; circles are the five candidate states; the judge decides per word how much each candidate counts. Dashed: reads the earlier words.}
\label{fig:inside}
\end{figure}

The sandbox works in four steps. Each is explained below twice: in plain words, and
in a technical explanation with the exact computation. Notation used throughout:
$h_t \in \mathbb{R}^{d}$ is the residual-stream state of token $t$ at the insertion
point ($d = 576$ for SmolLM2-135M), $x_t$ its normalised copy, $\Delta_{t,i}$ the change
proposed by reviewer $i$, $C_{t,i}$ the five candidate states, $p_{t,i}$ the judge's
weights; the token index $t$ is dropped where it is clear.

\subsubsection*{Step 1 --- four reviewers propose a corrected state}
\label{sec:roles}

\begin{plain}
Four reviewers each look at the word's state and propose a slightly different
version of it. They do \emph{not} do the same thing --- each has its own job and is
only allowed to look at certain things:
\begin{itemize}
\item \textbf{R1, the explorer}, looks only at the word itself and proposes \emph{another
reading} of it --- something the model may have under-weighted. Its rule: its
proposal must point in a new direction; it cannot just repeat or enlarge the original.
\item \textbf{R2, the checker}, also reads the earlier words of the question and
proposes the version that \emph{fits what was said before}.
\item \textbf{R4, the best-supported alternative}, reads the earlier words and the
proposals of R1 and R2, and combines them into the version \emph{with the most support}.
\item \textbf{R3, the doubter}, looks at how much R1, R2 and R4 \emph{disagree}. When
they disagree a lot, it proposes a version that keeps the question open instead of
committing too early.
\end{itemize}
At the start of training every reviewer proposes ``no change''; what they propose is
learned.
\end{plain}

\begin{tech}
\textbf{The generators.} Each reviewer $i$ owns a two-layer network
$F_i(z) = W_{2,i}\,\text{SiLU}(W_{1,i} z + b_{1,i}) + b_{2,i}$ with a hidden width of
$4d = 2304$. The output layer $W_{2,i}, b_{2,i}$ is initialised to zero, so every reviewer
proposes $\Delta_i = 0$ at the start of training and the model is exactly SmolLM2. Because
the gradient still reaches $W_{2,i}$ (its input $\text{SiLU}(\cdot)$ is not zero), the
reviewers start learning from the first step, in the same way as the zero-initialised
matrix of LoRA adapters.

\textbf{The text memory.} R2 and R4 may read the earlier tokens. Their keys are
$K = W_K\,\text{RMSNorm}(x_{\le t})$ and their values are the normalised states
$\text{RMSNorm}(x_{\le t})$ themselves, so the reviewers read \emph{the model's own
representation} of the context rather than a re-encoded copy. $\text{Attn}(q)$ denotes
multi-head attention (9 heads of 64 dimensions) of a query $q$ over this memory under a
causal mask; each reviewer has its own query matrix, so R2 and R4 can look for different
things in the same text.

\textbf{The four rules.}
\begin{align*}
\text{R1 (explore):}\quad & u = F_1(x),\qquad \Delta_1 = u - (u\cdot\hat x)\,\hat x, \qquad \hat x = x/\lVert x\rVert\\
\text{R2 (consistency):}\quad & \Delta_2 = F_2\big([\,x;\ \text{Attn}(W_{Q2}x)\,]\big)\\
\text{R4 (best-supported):}\quad & \Delta_4 = F_4\big([\,x;\ \text{Attn}(W_{Q4}x);\ \Delta_1;\ \Delta_2\,]\big)\\
\text{R3 (uncertainty):}\quad & \Delta_3 = F_3\big([\,x;\ \sigma\,]\big),\qquad \sigma = \sqrt{\text{Var}(\Delta_1,\Delta_2,\Delta_4) + \epsilon}
\end{align*}
$[\,\cdot\,;\,\cdot\,]$ is concatenation. R1's projection removes the component of its
proposal along $x$, so $\Delta_1 \perp x$ always holds: R1 can change the
\emph{direction} of the state (its reading) but can never just rescale or repeat it. R4
receives R1's and R2's proposals as inputs and is computed after them. R3's input
$\sigma$ is the per-channel standard deviation of the other three proposals, i.e.\ how
much they disagree in each of the 576 channels; $\epsilon = 10^{-6}$ keeps the gradient
finite at the start, where all proposals are equal. The candidates are
$C_0 = x$ (the model's own state, never changed) and $C_i = x + \Delta_i$ for
$i = 1, \dots, 4$.
\end{tech}

\begin{figure}[H]
\centering
\begin{tikzpicture}[scale=1.0]
  \draw[->,muted] (-0.3,0) -- (4.6,0); \draw[->,muted] (0,-0.3) -- (0,3.4);
  \fill[ink] (3.0,1.6) circle (2.6pt) node[below right,font=\small]{$x$ (model's state)};
  \draw[->,very thick,rA] (3.0,1.6) -- (2.2,3.1) node[above,font=\small]{$C_1 = x + \Delta_1$};
  \draw[dashed,muted] (0,0) -- (4.2,2.24);
  \draw[rA] (2.80,1.98) -- (3.17,2.18) -- (3.37,1.80);
  \node[lab,anchor=west] at (4.2,2.6) {$\Delta_1$ is at a right angle to $x$:\\R1 cannot simply stretch\\or copy the model's state};
\end{tikzpicture}
\caption{R1's rule, drawn in two of the 576 dimensions: whatever R1 learns, its change is orthogonal to the model's own state.}
\label{fig:r1}
\end{figure}

\subsubsection*{Step 2 --- the reviewers read each other once}

\begin{plain}
Before the judge decides, every reviewer sees the other proposals (and the original)
once and may adjust its own a little --- like a short discussion round. Each reviewer
adjusts in its own way, so they stay different. The original state takes part as a
reference but is never changed.
\end{plain}

\begin{tech}
\textbf{One communication round} (the number is configurable,
\code{sandbox\_rounds\_v01} $= 1$; each round has its own weights). All five candidates
are normalised, $n_i = \text{RMSNorm}(C_i)$, and every candidate reads the other four with
single-head attention, $\text{cross}_i = \sum_{j \ne i} a_{ij} V n_j$ with
$a_{i\cdot} = \text{softmax}_j\big((Q n_i)^\top (K n_j)/\sqrt d\big)$; the diagonal is
masked, so no candidate reads itself. Each reviewer then updates its own proposal:
\[
\Delta_i \leftarrow \Delta_i + \gamma_i \odot W_{o,i}\,\text{cross}_i,\qquad
\gamma_i = \sigma\big(W_{\gamma,i}[\,n_i;\ \text{cross}_i\,] + b_{\gamma,i}\big).
\]
$\gamma_i \in (0,1)^d$ is a per-channel \emph{repair gate}: it decides how much of what the
reviewer read it takes over. Its bias starts at $-2$ ($\gamma \approx 0.12$), so a reviewer
initially keeps its own position and only moves towards the others where training shows
that this helps. $W_{o,i}$ starts at zero (no change at step 0). The read is shared, but
$W_{\gamma,i}$ and $W_{o,i}$ belong to reviewer $i$, so each reacts to the discussion in its
own way and the roles do not merge. $C_0 = x$ is part of what is read but is never
updated: the original state stays available to the judge unchanged.
\end{tech}

\subsubsection*{Step 3 --- the judge decides, word by word}

\begin{plain}
A judge looks at all five versions --- the original and the four proposals --- together
with the word itself and the text so far, and gives each a weight. The weights add up
to 1. If no proposal helps, the judge can put almost all weight on the original, and
the word stays as SmolLM2 understood it. Because the judge decides per word, the
explorer may win at one word and the checker at the next. The weights are written
down during training, so we can see \emph{which} reviewer shaped a word --- a reason
for the choice, which the old average could never give.
\end{plain}

\begin{tech}
\textbf{Scoring.} The judge is a small network
$G = \text{Linear}(3d \to 256) \to \text{SiLU} \to \text{Linear}(256 \to 1)$ that scores
each candidate from three pieces of information:
\[
s_i = G\big([\,\text{RMSNorm}(C_i);\ x;\ \bar x_t\,]\big),\qquad i = 0,\dots,4,
\]
the candidate itself, the model's own state $x$ (what is being corrected), and
$\bar x_t$, the mean of the normalised states of tokens $1..t$ (a summary of the text so
far, computed causally). The same $G$ scores all five candidates, so the scores are
comparable.

\textbf{Weights.} $p = \text{softmax}(s_0,\dots,s_4)$, computed in 32-bit floating point
also under bf16 mixed precision, so the weights are always finite, non-negative and sum to
1 (tested). The last layer of $G$ starts at zero, so $p_i = 0.2$ for every candidate at
step 0. The decision is made independently for every token, so the dominant candidate can
change from word to word.

\textbf{Soft versus hard selection.} v0.1.0 uses \emph{soft} selection: the new state is
the $p$-weighted mixture (next step), which is differentiable, so the judge and all
reviewers receive gradients. \emph{Hard} selection (only the $\arg\max$ candidate) is not
differentiable and would need a straight-through or Gumbel-softmax estimator; it is
planned as a separate, controlled experiment. For analysis, \code{force\_candidate(i)}
replaces $p$ by a one-hot vector, which measures what each candidate alone does.
\end{tech}

\subsubsection*{Step 4 --- the improved state goes on}

\begin{plain}
The weighted mixture of the five versions becomes the word's new state, sized to fit
the model at that depth, and SmolLM2 continues with it.
\end{plain}

\begin{tech}
\textbf{Output.} The new state is the weighted mixture of the candidates, mapped back to
the scale of the stream:
\[
h'_t = h_t + \text{rms}(h_t)\sum_{i=1}^{4} p_{t,i}\,\Delta_{t,i}.
\]
Since $\sum_i p_{t,i} C_{t,i} = x_t + \sum_{i \ge 1} p_{t,i}\Delta_{t,i}$ (with $\Delta_0 = 0$),
this is exactly the spec's weighted selection $\sum_i p_i C_i$ in the normalised space,
added to the stream as a residual change. $\text{rms}(h_t)$ is treated as a constant
(no gradient), so the sandbox cannot learn to inflate the stream. The original state
always contributes with weight $p_0$, and when the judge puts all weight on $C_0$ the
token passes unchanged. With every $\Delta_i = 0$ at step 0, $h' = h$ bit for bit; the tests
check this after layer 1, after layer 2 and at the end, and also that switching the
sandbox off gives exactly SmolLM2.
\end{tech}

\subsection{An example of the intended behaviour}

\begin{plain}
\emph{Illustration of what the design is built for --- not a measured trace.}
Question: ``Is aspirin safe for children with a fever?'' At the word \emph{safe}, the
model's first reading leans towards ``a common painkiller, so yes''. R1 proposes the
other reading (``a risk''). R2 reads the earlier words \emph{children} and
\emph{fever} and proposes the version that fits them. R4 combines R1's alternative
with R2's context into the best-supported state. R1 and R2 disagree with the
original, so R3 proposes to keep the answer open. The judge gives most weight to R4
and R2 at this word, and the later layers continue from ``a risk for children''
instead of ``yes''. The sandbox can only use what SmolLM2 already knows --- it adds
no new medical facts --- but it gives the model a chance to question its first
reading before the error spreads to the rest of the answer.
\end{plain}

\subsection{How the reviewers evolve during training}

\begin{plain}
The reviewers are not fixed rules; they \textbf{learn while the model trains}. Two
things never change: \emph{what each reviewer may look at}, and R1's ``new
direction'' rule. Everything else --- what each one proposes, how they react to each
other, and how much the judge trusts each one --- is learned. After every batch,
each reviewer is nudged so that its proposal would have helped predict the next word,
\emph{in proportion to how much the judge used it}; the judge is nudged towards
reviewers that helped.

The weak point is a feedback loop: a reviewer the judge ignores gets almost no
nudge, stops improving, and is ignored even more (Figure~\ref{fig:loop}). That is
why every training log line shows how often each reviewer wins, how undecided the
judge is, and how strongly each reviewer is still learning.
\end{plain}

\begin{figure}[H]
\centering
\begin{tikzpicture}[node distance=0.9cm]
  \node[blk,minimum width=3.4cm] (a) {judge gives\\reviewer $i$ weight $p_i$};
  \node[blk,minimum width=3.4cm,right=1.6cm of a] (b) {reviewer $i$ learns\\at a rate $\propto p_i$};
  \node[blk,minimum width=3.4cm,below=1.0cm of $(a)!0.5!(b)$] (c) {reviewer $i$ helps\\more or less};
  \draw[flow] (a) -- (b); \draw[flow] (b) |- (c); \draw[flow] (c) -| (a);
  \node[blk,draw=lgreen,fill=plainbg,right=1.2cm of b,minimum width=4.2cm,align=left,font=\footnotesize] (fix) {\textbf{breaks the loop (optional):}\\candidate loss --- every reviewer\\learns even when ignored;\\early exploration --- the judge\\spreads its trust at first};
  \draw[->,thick,lgreen,dashed] (fix.west) -- (b.east);
\end{tikzpicture}
\caption{The training feedback loop that can make one reviewer dominate (spec \S8), and the two objectives built to break it.}
\label{fig:loop}
\end{figure}

\begin{tech}
\textbf{What is learned.} All parameters of the four generators, the two reviewer
queries, the memory key, the communication round and the judge are trained by
back-propagation together with SmolLM2 (or with SmolLM2 frozen during an optional
warm-up). Phase 1 uses only the next-token loss $\mathcal L$ on the answer tokens.

\textbf{How strongly each reviewer learns.} Differentiating the output equation gives
\begin{align*}
\frac{\partial\mathcal L}{\partial \Delta_i} &= p_i\,\text{rms}(h)\,\frac{\partial\mathcal L}{\partial h'}\quad (+\ \text{terms via R3, R4 and the round}),\\
\frac{\partial\mathcal L}{\partial s_i} &= p_i\,\text{rms}(h)\Big(\Delta_i - \textstyle\sum_j p_j\Delta_j\Big)\cdot\frac{\partial\mathcal L}{\partial h'} .
\end{align*}
The first line shows that reviewer $i$ is updated in proportion to its weight $p_i$: a
reviewer the judge uses gets a strong signal, an ignored one almost none. The second line
shows how the judge learns: it raises the score of a candidate whose change
$\Delta_i$ points further in the loss-reducing direction $-\partial\mathcal L/\partial h'$
than the current mixture $\sum_j p_j\Delta_j$, and lowers it otherwise. Together these
create the reinforcing loop of Figure~\ref{fig:loop}, the collapse risk named in spec \S8.

\textbf{What is monitored.} Every log line reports \code{sel} (the fraction of tokens on
which each of the five candidates has the largest weight), \code{H} (the mean entropy of
$p$; $\ln 5 = 1.609$ means undecided, $0$ means always the same candidate) and \code{gr}
(the gradient norm of each reviewer and of the judge). Each epoch additionally records
the mean weights, the $5\times5$ matrix of mean cosine similarities between candidates
(collapse $=$ off-diagonal values approaching 1) and the size of each $\Delta_i$.
\end{tech}

\subsection{Training each reviewer for its job (role objectives)}

\begin{plain}
On its own, training only says ``help predict the next word''; it does not tell R3
to be the doubter. Three extra training signals were built, each switched on
separately and \textbf{off by default}, so their effect can be measured one at a time:
\begin{enumerate}
\item \textbf{R3 learns uncertainty:} a tiny probe reads only R3's proposal and must
guess whether the model is about to get the next word right --- so R3's proposal has
to carry ``the model is unsure here''.
\item \textbf{No reviewer is starved:} each reviewer's own version is also scored on
the next word directly, even when the judge ignores it.
\item \textbf{The judge explores first:} for the first steps the judge is rewarded for
spreading its trust; the reward fades to zero.
\end{enumerate}
\end{plain}

\begin{tech}
The three objectives are added to the next-token loss with their own weights (all 0 by
default) and are computed only during training, on the answer tokens:
\[
\mathcal L_{total} = \mathcal L_{CE} + \lambda_{R3}\,\mathcal L_{R3} + \lambda_{cand}\,\mathcal L_{cand} - \lambda_t\,\bar H(p).
\]
\textbf{R3 correctness} (\code{--r3-loss} $\lambda_{R3}$). A linear probe $w^\top\Delta_3 + b$
(577 parameters, created only when the objective is on) reads R3's proposal and is trained
with binary cross-entropy against the target $\mathbb 1[\arg\max \text{logits}_t = y_{t+1}]$
--- whether the model's top prediction for the next token is correct. The target is
computed from the model's own output and detached (no gradient flows into it). Because
the probe sees only $\Delta_3$, the gradient forces R3's proposal to encode how likely the
model is to be wrong at this token --- a definition of ``uncertainty'' that can be checked.

\textbf{Candidate loss} (\code{--cand-loss} $\lambda_{cand}$).
$\mathcal L_{cand} = \tfrac14\sum_{i=1}^{4}\text{CE}\big(\text{head}(h + \text{rms}(h)\Delta_i),\ y\big)$:
each reviewer's \emph{own} candidate is passed through the output head and scored on the
next token, independent of $p_i$. A reviewer that the judge ignores still receives a full
learning signal, which breaks the loop described above. Scoring a candidate requires the
layers after the sandbox; at the end placement that is only the output head, so the
objective is limited to $k = -1$ (other depths would multiply the cost). The four
candidates are scored one after another to bound memory.

\textbf{Early exploration} (\code{--explore} $\lambda$, \code{--explore-steps} $N$).
$\bar H(p)$ is the mean entropy of the judge's weights over the answer tokens; subtracting
$\lambda_t\bar H(p)$ rewards spreading the weight over several candidates. The weight
decays linearly, $\lambda_t = \lambda\max(0, 1 - t/N)$, over the first $N$ optimiser steps of
stage 1 (default $N = 1000$), so it only prevents an early lock-in and does not force equal
use later (the spec explicitly rejects a fixed 20\,\% share per candidate).

Each objective is tested alone against the same run without it before any combination
(spec \S8.3); all three are logged per step (\code{r3=}, \code{cand=}, \code{entropy=}).
\end{tech}

\subsection{Size and cost}

\begin{center}
\begin{tabularx}{0.85\linewidth}{@{}Yr@{}}
\toprule
\textbf{Part (measured with \code{addon\_params()} on SmolLM2-135M)} & \textbf{Parameters}\\
\midrule
R1 explore & 2{,}657{,}088\\
R2 consistency (with its query) & 4{,}315{,}968\\
R3 uncertainty & 3{,}984{,}192\\
R4 best-supported (with its query) & 6{,}970{,}176\\
judge & 443{,}457\\
text memory + communication round & 5{,}314{,}176\\
norms of the insertion site & 1{,}728\\
\midrule
\textbf{sandbox v0.1.0} & \textbf{23{,}686{,}785}\\
sandbox v0.0.5, for comparison & 19{,}594{,}944\\
SmolLM2-135M backbone & 134{,}515{,}008\\
whole model v0.1.0 & 158{,}201{,}793\\
\bottomrule
\end{tabularx}
\end{center}

The sandbox is 20.9\,\% larger than v0.0.5 (within the $\pm$25\,\% agreed) and
17.6\,\% of the backbone. It runs once per token inside the normal forward pass ---
no second pass of SmolLM2, no generated drafts. The candidate loss is the only
part with a noticeable training cost (four extra output-layer evaluations on the
answer tokens), and it is optional.

\subsection{What changed compared with v0.0.5}

\begin{center}
\begin{tabularx}{\linewidth}{@{}p{3.0cm}YY@{}}
\toprule
 & \textbf{v0.0.5} & \textbf{v0.1.0}\\
\midrule
Starting points & $T_k = W_k h$, random orthogonal projections & $h + \Delta_i$, $\Delta_i = 0$ at start; R1 orthogonal by rule\\
Members & 4 identical structures & 4 roles with different inputs\\
Context & none (one word at a time) & R2 and R4 read the earlier words; the judge reads the text so far\\
Interaction & 4 debate rounds + 8 descent steps & 1 communication round (h is read, not changed)\\
Readout & mean of the four & learned judge over five, incl.\ the original\\
Output & $h + g(\text{RMSNorm}(c) - h)$, $g \approx 0.027$ & $h + \text{rms}(h)\sum p_i\Delta_i$\\
Placement & after the final norm & any depth (flag), study over six depths\\
Training per role & spread floor only & three optional role objectives\\
\bottomrule
\end{tabularx}
\end{center}

% =====================================================================
\section{The new evaluation: text generation}

\begin{plain}
Multiple-choice tests reward a model that leans towards one answer, which is
exactly what fooled us with v0.0.5. The new protocol asks the model to \emph{write}
answers. It first learns from examples that write out their reasoning step by step
before the answer, then from patient-style medical questions and answers, and is
finally judged on new questions: does its answer contain the key medical facts an
expert expects, and does it contradict any of them?
\end{plain}

\begin{figure}[H]
\centering
\resizebox{\linewidth}{!}{%
\begin{tikzpicture}[node distance=0.7cm]
  \node[blk,minimum width=3.6cm] (s1) {\textbf{stage 1}\\medical-o1-reasoning-SFT\\{\footnotesize question $\to$ reasoning $\to$ answer}};
  \node[blk,minimum width=3.0cm,right=of s1] (s2) {\textbf{stage 2}\\MedQuAD\\{\footnotesize fine-tune half}};
  \node[blk,minimum width=3.0cm,right=of s2] (sel) {checkpoint\\selection\\{\footnotesize MedQuAD validation}};
  \node[sbx,minimum width=3.6cm,right=of sel] (b) {\textbf{benchmark}\\MedLFQA\\{\footnotesize + MedQuAD bench half}};
  \draw[flow] (s1) -- (s2); \draw[flow] (s2) -- (sel); \draw[flow] (sel) -- (b);
  \node[blk,draw=lorange,fill=lorange!8,font=\footnotesize,below=0.7cm of s2,minimum width=6.4cm] (rm) {MedLFQA questions removed from stage 1 and MedQuAD\\before any split --- never trained on};
  \draw[->,thick,lorange,dashed] (b.south) |- (rm.east);
  \node[blk,font=\footnotesize,below=0.7cm of b,minimum width=3.6cm,align=left] (m) {perplexity of the reference\\ROUGE-L, token F1\\comprehensiveness, hallucination\\L2-off test, candidate forcing};
  \draw[flow] (b) -- (m);
\end{tikzpicture}}
\caption{The v0.1.0 protocol (\code{--stage1 reasoning --medlfqa}). Both arms use the same configuration file and flags; only \code{--variant} decides whether the sandbox is present.}
\label{fig:protocol}
\end{figure}

\begin{tech}
\begin{tabularx}{\linewidth}{@{}p{3.2cm}Y@{}}
\toprule
\textbf{Stage 1 data} & \code{FreedomIntelligence/medical-o1-reasoning-SFT} (en), 19{,}704 rows; target \code{<reasoning>\textbackslash nFinal answer: <response>}; 2{,}000 rows held out.\\
\textbf{Answer length} & SmolLM2 tokens of the targets, all 19{,}704 rows: median 535, p90 712, p99 935, max 1{,}472. A row longer than \code{max\_answer} is dropped, never cut (a cut row loses its final answer). 512 kept 8{,}164 rows (41.4\,\%); \textbf{768} (chosen) keeps 18{,}596 (94.4\,\%).\\
\textbf{Benchmark} & \code{dmis-lab/MedLFQA} (OLAPH, Jeong et al.\ 2024; LiveQA, MedicationQA, HealthSearchQA, K-QA): question, reference answer, \emph{must-have} and \emph{nice-to-have} statements.\\
\textbf{Factuality} & Following Manes et al.\ (2024): comprehensiveness = \% of must-have statements the answer \emph{entails}; hallucination = \% of all statements it \emph{contradicts}. Judge: \code{pritamdeka/PubMedBERT-MNLI-MedNLI}, used only for scoring, never trained; 1{,}000 seeded answers of up to 256 tokens.\\
\textbf{Statistics} & Paired bootstrap 95\,\% CI between the arms, resampling questions; a depth counts as better only if both the perplexity CI and a factuality CI exclude 0.\\
\bottomrule
\end{tabularx}
\end{tech}

% =====================================================================
\section{Results so far}

\begin{plain}
v0.1.0 has not been trained, because there is no GPU access at the moment. So
there is \textbf{no result yet on whether it reasons better}. What has been
established is that the new model and the new evaluation are correct and ready to
run: the code passes all automated tests, every part starts as an exact copy of
SmolLM2, and short runs on a laptop go through every stage without errors. These
checks also found and fixed one serious bug in the evaluation.
\end{plain}

\begin{center}
\begin{tabularx}{\linewidth}{@{}p{4.2cm}Y@{}}
\toprule
\textbf{Check} & \textbf{Result}\\
\midrule
Automated tests & 404 passed (346 before v0.1.0; 58 new for the sandbox v0.1.0, the role objectives, the data, the benchmark and the bug fix)\\
Start equals SmolLM2 & identical outputs at step 0 after layer 1, 2 and at the end; with the sandbox off always\\
Causality / cached generation & a later word never changes an earlier one; a cached step equals the full forward\\
Each reviewer learns & gradients reach all four roles and the judge (unit tests and every smoke-run log line)\\
CPU pipeline runs & four short runs (64 rows, 1 epoch each; after layer 20, at the end, baseline, end with all role objectives): no non-finite step; every stage and the MedLFQA benchmark ran. These are pipeline checks, \textbf{not results}\\
Size & 23{,}686{,}785 parameters, measured\\
\bottomrule
\end{tabularx}
\end{center}

\subsection{A bug the checks found (T65)}

\begin{plain}
In the first laptop runs the fact-checker gave a perfect score (100\,\% of the facts,
0\,\% contradictions) to answers of only 24 words --- impossible. The cause: the
fact-checker's word list had not been downloaded completely, and the software
quietly built an \emph{empty} one, so it read every word as ``unknown'' and said
``yes, this fact is in the answer'' to everything. The loader now refuses an
incomplete word list and downloads it properly. The same answers are now scored
47.2\,\% and 11.1\,\%, which is believable.
\end{plain}

\begin{tech}
\textbf{Cause.} The project loads models ``offline first'': it tries the local Hugging Face
cache and goes online only if that raises an error. For the judge, only
\code{config.json} was in the cache (it had been downloaded to read the label names). With
\code{transformers} 5, loading the tokenizer in this state does \emph{not} raise: it silently
builds a tokenizer with 5 entries, so every word became \code{[UNK]} and the NLI model
answered ``entailment'' for every pair. This was reproduced in a fresh, empty cache.

\textbf{Fix.} \code{load\_complete\_tokenizer} checks the vocabulary size after loading;
below a minimum (1{,}000 for SmolLM2, half the model's vocabulary for the judge) it makes
one online attempt and raises an error if the tokenizer is still incomplete --- an
incomplete tokenizer can no longer produce a silent result. Two tests cover the retry and
the error.

\textbf{Judge check.} Before relying on the judge, it was run on 12 hand-labelled medical
pairs (4 entailment, 5 neutral, 3 contradiction, including negations such as ``never
causes bleeding''): PubMedBERT-MNLI-MedNLI labelled 12/12 correctly; two general NLI
models (\code{cross-encoder/nli-deberta-v3-base},
\code{MoritzLaurer/DeBERTa-v3-base-mnli-fever-anli}) 8/12 each. Twelve pairs are a
sanity check, not a validation; the judge's scores are only compared between arms under
identical conditions.
\end{tech}

% =====================================================================
\section{The experiment plan}

\begin{center}
\begin{tabularx}{\linewidth}{@{}p{2.6cm}Yp{4.4cm}@{}}
\toprule
\textbf{Phase} & \textbf{What} & \textbf{Decision rule}\\
\midrule
3 sanity & one seed; baseline, sandbox at the end and after layer 20; SmolLM2 frozen for the first 2{,}000 steps & no non-finite steps; every role gets gradient; judge not stuck on one candidate; candidates not identical\\
4 placement & one sandbox per run after layers 5, 10, 15, 20, 25 or at the end; same seed, data and flags; no warm-up (as the baseline); then three seeds for the best two & perplexity CI and a factuality CI against the baseline both exclude 0\\
5 role objectives & at the best depth: R3 correctness, candidate loss, early exploration --- each alone against the same run without it & the objective improves factuality or prevents collapse, with a CI that excludes 0\\
6 later & hard selection, more rounds, several insertion points & only if 4--5 show a gain\\
\bottomrule
\end{tabularx}
\end{center}

Every run ends with the L2-off test (the same model with the sandbox switched off)
and the candidate test (\code{gen\_eval candidates}: each of the five candidates
forced in turn). Together they show whether the model actually depends on the
sandbox and which reviewer carries the change. Commands: \code{run.md} A3t.

% =====================================================================
\section{Risks and open questions}

\begin{plain}
\begin{itemize}
\item \textbf{The reviewers may not become what their names say.} Their inputs and
rules are fixed, but what they learn depends on training. Whether R1 really explores
and R2 really checks is a result to measure, not a promise.
\item \textbf{One reviewer may take over} (the feedback loop above). The role
objectives exist to prevent this, and the logs show it early.
\item \textbf{The judge may learn slowly.} In the laptop runs its learning signal was
far smaller than the reviewers', and it stayed undecided over 14--22 steps. That is
expected for such short runs; if it persists in Phase 3, the judge gets its own,
higher learning rate (as the old gate had).
\item \textbf{No new knowledge.} The sandbox can only use what SmolLM2 knows better.
\item \textbf{The fact-checker is a model too.} It passed a small hand check; its
scores are compared between arms under identical conditions, not read as absolute truth.
\end{itemize}
\end{plain}

% =====================================================================
\section{Summary}

\begin{keybox}[title={Summary}]
\begin{itemize}[leftmargin=1.2em]
\item The first sandbox did not improve reasoning; its gain was a ``yes'' preference.
\item \textbf{Ekdant v0.1.0} rebuilds it as a team of four role agents inside the model ---
explore, check consistency, express uncertainty, build the best-supported alternative ---
plus the model's own state and a learned judge that decides per word and records why.
\item Each reviewer learns during training; three optional objectives push the roles
towards their jobs and protect against one reviewer taking over.
\item The benchmark is now written answers checked for medical facts (MedLFQA).
\item Built, tested and checked end-to-end; \textbf{training starts when a GPU is available}
--- first the sanity run, then the placement study, then the role objectives one at a time.
\end{itemize}
\end{keybox}

% =====================================================================
\appendix
\section{Glossary}

\begin{center}
\begin{tabularx}{\linewidth}{@{}>{\raggedright\arraybackslash}p{4.2cm}Y@{}}
\toprule
State ($h$) & the 576 numbers SmolLM2 holds for a word at a given layer\\
Candidate & one of five versions of a word's state: the original ($C_0$) or a reviewer's proposal ($C_1..C_4$)\\
Role generator / reviewer & a small network that proposes a change $\Delta_i$ to the state, with its own inputs and rule\\
Judge / selection gate & a small network that scores the five candidates per word; softmax turns scores into weights\\
Soft selection & the new state is the weighted mix of the candidates (trainable directly)\\
Insertion point & the layer after which the sandbox runs (\code{--addon-layer})\\
Collapse & one candidate wins everywhere or all candidates become the same\\
Entailment / contradiction & the answer implies / denies a statement (judged by the NLI model)\\
Comprehensiveness & \% of the must-have statements an answer entails\\
Hallucination & \% of all statements an answer contradicts\\
Perplexity & how surprised the model is by the reference answer (lower is better)\\
L2-off test & the same trained model with the sandbox switched off\\
\bottomrule
\end{tabularx}
\end{center}

\section{Reproducibility}

\begin{center}\small
\begin{tabularx}{\linewidth}{@{}>{\raggedright\arraybackslash}p{3.6cm}Y@{}}
\toprule
Specification and explanation & \code{04-Model-DeepDive/30\_L2\_v6\_Roles\_Spec.md}\\
Implementation plan & \code{04-Model-DeepDive/30a\_L2\_v6\_Roles\_Plan.md}\\
Code & \code{09-Code/src/ekdant/sandbox.py} (\code{SandboxV01}), \code{model.py}, \code{train\_gen.py}, \code{gen\_data.py}, \code{gen\_eval.py}\\
Configuration (both arms) & \code{09-Code/configs/ekdant\_v0.1.0.json}\\
Tests & \code{09-Code/tests/test\_sandbox\_v01.py} (43) and additions in \code{test\_gen\_data.py}, \code{test\_gen\_eval.py}, \code{test\_backbone.py}\\
Run commands & \code{run.md} A3t; issues T64, T65 in \code{07-Management/10\_Known\_Issues.md}\\
\bottomrule
\end{tabularx}
\end{center}

\end{document}
